% \begin{document}
\chapter{Discrete Probability (a short reminder)}

Un breve ripasso di probabilità discreta: serve per ragionare sulla sicurezza dei cifrari (es. la perfect secrecy dell'OTP e l'uniformità del keystream).

\section{Distribuzione di probabilità}

\dfn{Universe / Sample space}{
  $U$ è un \textbf{insieme finito}, detto \textit{Universe} o \textit{Sample space}.
}
Esempi: lancio di una moneta $U = \{\text{heads}, \text{tail}\}$ (oppure $U = \{0,1\}$); lancio di un dado $U = \{1,2,3,4,5,6\}$. In crittografia spesso $U = \{0,1\}^n$.

\dfn{Distribuzione di probabilità}{
  Una \textit{distribuzione di probabilità} $P$ su $U$ è una funzione
  \[
    P : U \to [0,1] \qquad \text{tale che} \qquad \sum_{x \in U} P(x) = 1
  \]
}
Esempi: moneta $P(\text{heads}) = P(\text{tail}) = 1/2$; dado $P(1) = \dots = P(6) = 1/6$. Oppure su $U = \{0,1\}^2 = \{00,01,10,11\}$: $P(00) = 1/2$, $P(01) = 1/8$, $P(10) = 1/4$, $P(11) = 1/8$ (somma $= 1$).

\subsection{Distribuzioni particolari}
\begin{enumerate}
  \item \textbf{Distribuzione uniforme}: $\forall x \in U,\ P(x) = 1/|U|$.
  \item \textbf{Point distribution} in $x_0$: $P(x_0) = 1$ e $\forall x \neq x_0,\ P(x) = 0$.
\end{enumerate}
\dots e molte altre.

\section{Eventi}

Sia $U$ un universe con distribuzione $P$.
\dfn{Evento}{
  Un \textit{evento} è un sottoinsieme $A \subseteq U$. La sua probabilità è
  \[
    \Pr[A] = \sum_{x \in A} P(x)
  \]
}
\nt{ Si noti che $\Pr[U] = 1$. }

\ex{Dado}{
  $U = \{1,\dots,6\}$ uniforme, $A = \{1,3,5\}$ (numeri dispari):
  \[
    \Pr[A] = \tfrac16 + \tfrac16 + \tfrac16 = \tfrac12
  \]
}

\ex{Bit meno significativi}{
  $U = \{0,1\}^8$ con distribuzione uniforme ($P(x) = 1/2^8$ per ogni $x$). Sia
  \[
    A = \{\, x \in U : \mathrm{lsb}_2(x) = 11 \,\} \subseteq U
  \]
  (gli $x$ i cui due bit meno significativi sono \texttt{11}, cioè della forma \texttt{\_\_\_\_\_\_11}). Poiché $\Pr[A] = \frac{1}{2^8}\cdot|A|$ e $|A| = 2^6$ (i 6 bit liberi):
  \[
    \Pr[A] = \frac{2^6}{2^8} = \frac14
  \]
}

\subsection{Unione di eventi}
Dati due eventi $A_1, A_2$, anche $A_1 \cup A_2$ è un evento e:
\[
  \Pr[A_1 \cup A_2] = \Pr[A_1] + \Pr[A_2] - \Pr[A_1 \cap A_2]
\]
\mprop{Union bound}{
  \[
    \Pr[A_1 \cup A_2] \le \Pr[A_1] + \Pr[A_2]
  \]
}
Se inoltre $A_1 \cap A_2 = \varnothing$ (eventi disgiunti) allora vale l'uguaglianza:
\[
  A_1 \cap A_2 = \varnothing \implies \Pr[A_1 \cup A_2] = \Pr[A_1] + \Pr[A_2]
\]

\section{Variabili aleatorie}

\dfn{Variabile aleatoria}{
  Una \textit{random variable} $X$ è una funzione $X : U \to V$. $X$ assume valori in $V$ e \textit{induce} una distribuzione su $V$.
}

\ex{Dado pari/dispari}{
  $U = \{1,\dots,6\}$ uniforme, $X : U \to V = \{\text{``even''}, \text{``odd''}\}$ con
  \[
    X(2) = X(4) = X(6) = \text{``even''}, \qquad X(1) = X(3) = X(5) = \text{``odd''}
  \]
  Allora $\Pr[X = \text{``even''}] = \tfrac12$ e $\Pr[X = \text{``odd''}] = \tfrac12$.
}

\subsection{La variabile aleatoria uniforme}
Sia $S$ un insieme (es. $S = \{0,1\}^n$). Scriviamo
\[
  X \leftarrow S
\]
per indicare che $X$ è una \textbf{variabile aleatoria uniforme} su $S$, cioè
\[
  \forall a \in S : \quad \Pr[X = a] = \frac{1}{|S|}
\]

\subsection{Algoritmi randomizzati}
\begin{itemize}
  \item \textbf{Algoritmo deterministico}: $y \leftarrow A(m)$, a ogni input corrisponde un unico output.
  \item \textbf{Algoritmo randomizzato}: l'output è una variabile aleatoria $y \leftarrow A(m)$ (lo stesso input può dare output diversi, secondo una distribuzione).
\end{itemize}

\section{Indipendenza}

\dfn{Eventi indipendenti}{
  Due eventi $A$ e $B$ sono \textit{indipendenti} se
  \[
    \Pr[A \cap B] = \Pr[A] \cdot \Pr[B]
  \]
}
\dfn{Variabili aleatorie indipendenti}{
  Due variabili aleatorie $X, Y$ a valori in $V$ sono \textit{indipendenti} se
  \[
    \forall a, b \in V : \quad \Pr[X = a \text{ and } Y = b] = \Pr[X = a] \cdot \Pr[Y = b]
  \]
}

\section{XOR e la sua proprietà fondamentale}

Lo XOR di due stringhe in $\{0,1\}^n$ è la loro somma bit a bit modulo 2 ($\oplus$); $\mathrm{XOR}(X)$ indica lo xor di tutti i bit della stringa $X$ (la parità). La proprietà che lo rende utilissimo in crittografia:

\thm{Proprietà fondamentale dello XOR}{
  Siano:
  \begin{enumerate}
    \item $X$ una variabile aleatoria su $\{0,1\}^n$ con distribuzione \textbf{uniforme};
    \item $Y$ una variabile aleatoria su $\{0,1\}^n$ con distribuzione \textbf{arbitraria};
    \item $X$ e $Y$ \textbf{indipendenti}.
  \end{enumerate}
  Allora $Z := Y \oplus X$ è una variabile aleatoria \textbf{uniforme} su $\{0,1\}^n$.
}

\pf{Dimostrazione (per $n = 1$)}{
  Sia $\Pr[Y=0] = p_0$ e $\Pr[Y=1] = p_1$ (con $p_0 + p_1 = 1$), mentre $X$ è uniforme: $\Pr[X=0] = \Pr[X=1] = \tfrac12$. Per indipendenza, la coppia $(X,Y)$ ha probabilità $\Pr[X=x]\cdot\Pr[Y=y]$. Allora $Z = 0$ accade quando $X = Y$:
  \[
    \Pr[Z = 0] = \Pr[(X,Y) = (0,0)] + \Pr[(X,Y) = (1,1)] = \frac{p_0}{2} + \frac{p_1}{2} = \frac{p_0 + p_1}{2} = \frac12
  \]
  Di conseguenza $\Pr[Z = 1] = \tfrac12$, quindi $Z$ è uniforme.
}
\nt{
  È esattamente questa proprietà che dà la perfect secrecy all'OTP: facendo lo XOR del messaggio (distribuzione \emph{qualsiasi}) con una chiave \emph{uniforme} e indipendente, il ciphertext risulta uniforme e non rivela informazione sul plaintext.
}

\section{Il paradosso del compleanno (birthday paradox)}

Siano $r_1, \dots, r_n \in U$ variabili aleatorie \textbf{independent identically distributed} (i.i.d.).
\thm{Birthday paradox}{
  Quando $n = 1.2 \times |U|^{1/2}$, allora
  \[
    \Pr[\exists\, i \neq j : r_i = r_j] \ge \frac12
  \]
  cioè con circa $\sqrt{|U|}$ campioni si ha una \textbf{collisione} con probabilità $\ge 1/2$.
}

\ex{Compleanni}{
  $U = \{1,2,\dots,366\}$. Con $n = 1.2 \times \sqrt{366} \approx 23$ persone, due hanno lo stesso compleanno con probabilità $\ge 1/2$ (sorprendentemente poche!).
}
\ex{Spazio crittografico}{
  $U = \{0,1\}^{128}$. Dopo aver campionato circa $2^{64}$ messaggi random da $U$, è probabile che due coincidano. Questo è il motivo per cui un valore casuale a $n$ bit "regge" solo fino a $\sim 2^{n/2}$ usi.
}
\nt{
  Vedremo più avanti un esempio pratico con \textbf{WEP} (Wired Equivalent Privacy, IEEE 802.11): l'IV a 24 bit collide dopo $\sim 2^{12}$ frame per il birthday paradox (in realtà si ripete del tutto dopo $2^{24}$). La probabilità di collisione cresce ad "S" con il numero di campioni $n$ e raggiunge $1/2$ attorno a $1.2\sqrt{|U|}$.
}
% \end{document}
